\documentclass[12pt]{article}
\usepackage[margin=1in]{geometry}
\usepackage{hyperref}
\hypersetup{colorlinks=true, linkcolor=blue, urlcolor=blue, citecolor=blue}

\title{\textbf{A Complexity-Theoretic Foundation for Metrology \\ \large Book II: The Attractor-Basin Program for Type B Certification}}
\author{Dmytro Surdu}
\date{\today}

\begin{document}
\maketitle

\begin{abstract}
\noindent
This second volume extends the complexity-theoretic framework of Book~I into the domain of \textbf{Type B measurements}—where knowledge arises from models, constraints, and expert judgement rather than statistical sampling. We present the \emph{Attractor Basin Equivalence Theorem (ABET)} as the exact structural criterion for deterministic certifiability of tail properties, unifying the epistemic notion of a short proof with the physical reality of a reachable, invariant basin. Along the way, we establish the \emph{Minimality Theorem} for basin invariance, the per-step \emph{Double-Bound} method for certifying belief-state updates, and the \emph{Minimal Certification Standard} that forces the unique Observation--Adaptation--Restitution anatomy of any certifiable process. This book offers a rigorous bridge between the formal world of complexity theory and the practical challenges of designing measurement systems that meet the strongest possible certification standards.
\end{abstract}

\section*{Foreword and Guide to This Collection}

The first volume built a new foundation for \textbf{Type A} measurement---statistical sampling---by showing that deterministic certification demands a finite-state, information-losing structure bounded by guard-bands. This second book continues the story in a different landscape: \textbf{Type B} measurements, where the raw randomness of repeated sampling is replaced by the persistent structure of a model.

Here, the key question is: \emph{What physical and algorithmic structure makes a Type B claim certifiable in the same strong sense as in Book~I?} The answer leads us to \emph{basins}---regions of belief-state space that, once entered, guarantee future safety---and to a tight correspondence between such basins and certifiable tail predicates.

As in the first volume, the path to this result is neither linear nor purely theoretical. We begin with a concrete case study, abstract its key components into a formal model, and only then push the theory to its necessary and sufficient conditions. The narrative alternates between \textbf{constructive examples} and \textbf{formal equivalences}, ensuring that each definition earns its keep in proving a major theorem. On its way, developing theory gives rise to the most widely accepted notions and objects in metrology.

The reader is encouraged to follow the chapters in the order presented. Part~I develops per-step certification and its complexity bounds. Part~II establishes the structural machinery of basins and proves the central ABET equivalence. Part~III connects Type B results back to Type A and analyzes the anatomy of a certifiable measurement in its simplest canonical form.

\newpage
\tableofcontents
\newpage

\part[From Examples to Per-Step Certification]{From Examples to Per-Step Certification}

\noindent
This first part lays the groundwork for the Attractor-Basin Program. We start with the motivating problem: how can we hope to deterministically certify the long-term behaviour of a system driven by structured, non-statistical information? The journey begins with a realistic Type~B case study and proceeds to a formal belief-state model. We then prove that per-step testability is equivalent to controlled information loss, and develop the \emph{Double-Bound} method that will anchor the rest of the book.

\section*{Chapter 1: The Quest for Deterministic Certification in Metrology}
Introduces the challenge of $\delta \to 0$ certainty for Type~B claims, the $\Pi^0_2$ barrier, and the roadmap: per-step certification first, tail certification via structural constraints.

\section*{Chapter 2: The Outlier-Robust Kalman Filter---A Primer on Type B Certification}
A worked example of a belief-state system with robust, polynomial-time updates that lose information each step. Serves as a touchstone for the general theory.

\section*{Chapter 3: The Formalism of Belief-State Metrology}
Defines the belief-state $\Theta_t$, its compact state space, and the deterministic \emph{Emit/Update} cycle. Introduces multiple equivalent notions of per-step information loss.

\section*{Chapter 4: Per-Step Testability and its Equivalence Proof}
Shows that per-step deterministic certification is equivalent to controlled information loss. Presents the \emph{Double-Bound} method: a numerical contraction bound plus a refinement procedure ensuring complexity feasibility. Introduces \emph{epistemic losslessness} for basin interiors.

\part[Basins and Tail Certification]{Basins and Tail Certification}

\noindent
With per-step certification in hand, we turn to the structural problem: how to make claims about the infinite tail of a measurement sequence. The central objects here are \emph{basins}—regions of the belief-state space that guarantee safety once entered. We establish the minimal geometric condition that makes basins invariant, and then prove the \emph{Attractor Basin Equivalence Theorem}, which characterises exactly when a tail predicate is certifiable.

\section*{Chapter 5: The Minimal Condition for Invariant Basins---Local Contraction}
Axiomatizes basins and proves the \emph{Minimality Theorem}: uniform strict local contraction is both necessary and sufficient for invariance under mild conditions, making it the minimal geometric regularity for tail safety.

\section*{Chapter 6: The Attractor Basin Equivalence Theorem (ABET)}
The main result: a tail predicate is deterministically certifiable \emph{iff} there exists a reachable, invariant basin meeting emission safety and verifiability conditions. Establishes the exact bridge between logical certifiability and physical structure.

\part[Unification and Anatomy]{Unification and Anatomy}

\noindent
The final part unites the Type~A and Type~B worlds and examines the internal anatomy of a certifiable measurement. We show how Type~B dynamics can be cast as guard-band predicates recognisable in the Type~A framework, and how the \emph{Minimal Certification Standard} forces a unique decomposition of any certifiable process into its three canonical aspects. The \emph{Simple Test} emerges as the maximally compressed, canonical form of certifiable behaviour.

\section*{Chapter 7: Unification of Type A and Type B Measurements}
Introduces the Long-Trajectory Cast to express Type~B dynamics as guard-band predicates, unifying both measurement types in the same logical framework.

\section*{Chapter 8: Analyzing the Testable Measurement}
Defines the \emph{Minimal Certification Standard} $\{\varepsilon_n\}$, which induces equivalent state, input, and output guard-bands. Proves the unique decomposition into \textbf{Observation}, \textbf{Adaptation}, and \textbf{Restitution}. Shows that, in the \emph{Simple Test}, all post-entry certified outputs collapse to a single epistemic tick.

\vspace{1em}
\noindent\textit{Welcome to the Attractor-Basin Program---the second step in a unified theory of certifiable measurement.}

\end{document}
